\section{Open problems}\label{sec:status}

Six problems are open. Heat and an
applied tide fix a body's interior in principle, with no null space and through the
same logarithm as the static field (Section~\ref{sec:second-order}); how much of a real
body, the Earth for one, a finite record at finite sensitivity over a finite
band actually recovers has not been worked out. The pair of arrangements that
gravity cannot tell apart at one time, Eq.~\eqref{eq:dp-energy}, is a test of gravitational
collapse that is described here and not designed. The reading of a body by
an impulsive tide and a weighing, Eq.~\eqref{eq:push}, is exact in Newton's theory with
the weight of energy added, and for an impulsive plane wave it is
Eq.~\eqref{eq:shock}. A particle without mass that passes,
Eq.~\eqref{eq:massless-pass}, is the simplest wave sent in from infinity; the
statement for a general wave, with the mass read at infinity for the weight,
is a problem of scattering that has not been solved. The brackets of the
moments are given here to the first order in $1/c^2$,
Eq.~\eqref{eq:moment-bracket}; whether the moments of the full theory close
on a known algebra is open \cite{CompereOliveriSeraj2018}. The ceiling on
sensitivity of Section~\ref{sec:near-field} is a heuristic estimate; a bound
needs the coupling, the resources and the task to be stated exactly. And the
friction of real materials at the frequencies of heat, on which every thermal
count in this paper rests, deserves a microscopic calculation for at least one
crystal, phonon by phonon, of the kind already made for absorption by crystal
detectors \cite{KahnSchutteEngelTrickle2024}, and for the deep Earth at its
temperatures and pressures.
